% Part: incompleteness
% Chapter: representability-in-q
% Section: composition-representable

\documentclass[../../../include/open-logic-section]{subfiles}

\begin{document}

\olfileid{inc}{req}{cmp}
\olsection{মিশ্রণ $\Th{Q}$-তে উপস্থাপনযোগ্য}

ধরি $h$-কে এভাবে সংজ্ঞায়িত করা হয়েছে:
\[
h(x_0,\dots,x_{l-1}) = f(g_0(x_0,\dots,x_{l-1}), \dots,
g_{k-1}(x_0,\dots,x_{l-1})).
\]
এখানে $f$, $g_0$, \dots,~$g_{k-1}$ অপেক্ষকগুলিকে যথাক্রমে উপস্থাপন করে
এমন !!{formula}s $!A_f, !A_{g_0}, \dots, !A_{g_{k-1}}$ আমরা ইতিমধ্যে
পেয়েছি। $h$-কে উপস্থাপন করে এমন একটি !!a{formula}~$!A_h$ খুঁজতে হবে।

প্রথমে সব অপেক্ষক $1$-স্থানীয় এমন একটি সরল ক্ষেত্র নিই; অর্থাৎ
$h(x) = f(g(x))$ বিবেচনা করি। $!A_f(y, z)$ যদি $f$-কে এবং $!A_g(x, y)$
যদি $g$-কে উপস্থাপন করে, তাহলে $h$-কে উপস্থাপনকারী
!!a{formula}~$!A_h(x, z)$ দরকার। লক্ষ করো, $h(x) = z$ যদি এবং কেবল যদি এমন
একটি~$y$ থাকে যাতে $z = f(y)$ ও $y = g(x)$ উভয়ই হয়। ($h(x) = z$ হলে
$g(x)$ এমন একটি~$y$; আর এমন $y$ থাকলে $y = g(x)$ ও $z = f(y)$ থেকে
$z = f(g(x))$।) এতে বোঝা যায়,
$\lexists[y][(!A_g(x, y) \land !A_f(y, z))]$ হলো $!A_h(x, z)$-এর ভালো
প্রার্থী। শুধু যাচাই করতে হবে যে $\Th{Q}$ প্রাসঙ্গিক !!{formula}s প্রমাণ
করে।

\begin{prop}
\ollabel{prop:rep1}
$h(n) = m$ হলে $\Th{Q} \Proves !A_h(\num{n}, \num{m})$।
\end{prop}

\begin{proof}
ধরি $h(n) = m$, অর্থাৎ $f(g(n)) = m$। $k = g(n)$ ধরি। তাহলে
\begin{align*}
  \Th{Q} & \Proves !A_g(\num{n}, \num{k})
  \intertext{কারণ $!A_g$ অপেক্ষক~$g$-কে উপস্থাপন করে, এবং}
  \Th{Q} & \Proves !A_f(\num{k}, \num{m})
  \intertext{কারণ $!A_f$ অপেক্ষক~$f$-কে উপস্থাপন করে। তাই}
  \Th{Q} & \Proves !A_g(\num{n}, \num{k}) \land !A_f(\num{k}, \num{m})
  \intertext{এবং ফলত}
  \Th{Q} & \Proves \lexists[y][(!A_g(\num{n}, y) \land !A_f(y, \num{m}))],
\end{align*}
অর্থাৎ $\Th{Q} \Proves !A_h(\num{n}, \num{m})$।
\end{proof}

\begin{prop}
\ollabel{prop:rep2}
$h(n) = m$ হলে $\Th{Q} \Proves \lforall[z][(!A_h(\num{n}, z) \lif
  z = \num{m})]$।
\end{prop}

\begin{proof}
ধরি $h(n) = m$, অর্থাৎ $f(g(n)) = m$। $k = g(n)$ ধরি। তাহলে
\begin{align*}
  \Th{Q} & \Proves \lforall[y][(!A_g(\num{n}, y) \lif \eq[y][\num{k}])]
  \intertext{কারণ $!A_g$ অপেক্ষক~$g$-কে উপস্থাপন করে, এবং}
  \Th{Q} & \Proves \lforall[z][(!A_f(\num{k}, z) \lif \eq[z][\num{m}])]
  \intertext{কারণ $!A_f$ অপেক্ষক~$f$-কে উপস্থাপন করে। অল্প একটু যুক্তিবিদ্যা
    ব্যবহার করে আরও দেখানো যায় যে}
  \Th{Q} & \Proves \lforall[z][(\lexists[y][(!A_g(\num{n}, y) \land
      !A_f(y, z))] \lif \eq[z][\num{m}])].
\end{align*}
অর্থাৎ $\Th{Q} \Proves \lforall[y][(!A_h(\num n, y) \lif \eq[y][\num m])]$।
\end{proof}

$f$ ও~$g_i$-এর স্থানসংখ্যা $1$-এর বেশি হলেও একই ধারণা কাজে লাগে।

\begin{prop}
\ollabel{prop:rep-composition}
$!A_f(y_0, \dots, y_{k-1}, z)$ যদি $\Th{Q}$-তে
$f(y_0, \dots, y_{k-1})$-কে এবং
$!A_{g_i}(x_0, \dots, x_{l-1}, y)$ যদি $\Th{Q}$-তে
$g_i(x_0, \dots, x_{l-1})$-কে উপস্থাপন করে, তাহলে
\begin{multline*}
  \lexists[y_0][\dots \lexists[y_{k-1}][(!A_{g_0}(x_0,\dots,x_{l-1},y_0) \land
      \dots \land {}\\
  !A_{g_{k-1}}(x_0,\dots,x_{l-1},y_{k-1}) \land !A_f(y_0,\dots,y_{k-1},z))]]
\end{multline*}
উপস্থাপন করে
\[
h(x_0, \dots, x_{l-1}) = f(g_0(x_0, \dots, x_{l-1}), \dots, g_{k-1}(x_0,
\dots, x_{l-1})).
\]
% BN-SRC-222 TARGET-CORRECTION-BEGIN
\par\smallskip\noindent\textnormal{\small\textbf{উৎস-সংশোধন (BN-SRC-222):} হিমায়িত উপস্থাপক সূত্রে প্রথম অস্তিত্বমূলক পরিমাণসূচকের চলরাশি-তালিকায় $\dots$ ঢুকে গিয়েছিল, অকালেই দুটি বন্ধনী বন্ধ হয়েছিল, এবং শেষের পরিমাণসূচকগুলির পরিসর খোলা ছিল। বাংলা পাঠে একই উদ্দেশ্যপ্রকাশক সূত্রকে $y_0$ থেকে $y_{k-1}$ পর্যন্ত সঠিকভাবে বদ্ধ ও সুষম বন্ধনীতে লেখা হয়েছে।} % BN-SRC-222 TARGET-CORRECTION-END
\end{prop}

\begin{proof}
অনুশীলন।
\end{proof}

\begin{prob}
\olref[inc][req][cmp]{prop:rep1} ও
\olref[inc][req][cmp]{prop:rep2}-এর প্রমাণকে নির্দেশক হিসেবে ব্যবহার করে
\olref[inc][req][cmp]{prop:rep-composition}-এর প্রমাণ বিশদে সম্পন্ন করো।
% BN-SRC-223 TARGET-CORRECTION-BEGIN
\par\smallskip\noindent\textnormal{\small\textbf{উৎস-সংশোধন (BN-SRC-223):} হিমায়িত সমস্যায় নির্দেশক হিসেবে $\mathrm{prop{:}rep2}$-কে দুবার উল্লেখ করা হয়েছিল। বাংলা পাঠে আগের দুটি পৃথক প্রস্তাবের সঠিক জোড়া $\mathrm{prop{:}rep1}$ ও $\mathrm{prop{:}rep2}$ রাখা হয়েছে।} % BN-SRC-223 TARGET-CORRECTION-END
\end{prob}

\end{document}
